\documentclass[11pt,a4paper]{article}
\usepackage[utf8]{inputenc}
\usepackage[T1]{fontenc}
\usepackage[brazilian]{babel}
\usepackage{geometry}
\usepackage{graphicx}
\usepackage{booktabs}
\usepackage{float}
\usepackage{enumitem}
\usepackage{amsmath}
\usepackage{xcolor}
\usepackage{tcolorbox}
\geometry{margin=2.5cm}
\setlength{\parindent}{0pt}
\setlength{\parskip}{0.5em}
\tcbuselibrary{skins,breakable}
\newtcolorbox{nota}[1][]{colback=yellow!10,colframe=orange!60!black,fonttitle=\bfseries,title=#1,breakable}

\begin{document}
\begin{center}
  
\includegraphics[width=0.16\textwidth]{../../../figures_global/unb.png}\\[0.3cm]
  {\large\textbf{GABARITO DO PROFESSOR -- USO RESTRITO}}\\[0.2cm]
  {\large FGA0092 -- Princípios de Comunicação para Engenharia}\\[0.2cm]
  \rule{12cm}{0.8pt}\\[0.2cm]
  {\Large\textbf{Prática 02: Modulação AM DSB-SC}}\\[0.2cm]
  \rule{12cm}{0.8pt}
\end{center}

\section*{Gráfico esperado}
\begin{center}
  \includegraphics[width=0.88\textwidth]{grafico_esperado.png}
\end{center}
\textit{O gráfico acima foi gerado pelo script \texttt{gerar\_grafico.py} desta pasta.}

\section*{Respostas esperadas}

\begin{enumerate}[label=\textbf{Q\arabic*.}, leftmargin=*, itemsep=1em]

  \item \textbf{Por que o espectro do DSB-SC não apresenta pico na frequência da portadora?}

    \textbf{Resposta esperada:} Para $s(t) = m(t)\cos(2\pi f_c t)$ com
    $m(t) = A_m\cos(2\pi f_m t)$, o produto resulta em
    $S(f) = \tfrac{A_m}{2}[\delta(f-(f_c-f_m)) + \delta(f-(f_c+f_m))]$.
    O coeficiente em $f_c$ é zero porque não há componente constante (DC) em $m(t)$.
    No espectro do GNU Radio, a ausência do pico em $f_c = 10000$\,Hz confirma que
    toda a energia está nas bandas laterais.

    \begin{nota}[Erros comuns]
      Não relacionar a ausência da portadora com a falta de componente DC na mensagem.
      A explicação matemática usando o produto de cossenos deve ser apresentada.
    \end{nota}

  \item \textbf{Como a variação de $f_m$ afeta as posições das bandas laterais no espectro?}

    \textbf{Resposta esperada:} As bandas laterais são $f_\text{BL} = f_c \pm f_m$.
    Para $f_c = 10000$\,Hz: $f_m=500 \to 9500/10500$\,Hz; $f_m=1000 \to 9000/11000$\,Hz;
    $f_m=2000 \to 8000/12000$\,Hz; $f_m=4000 \to 6000/14000$\,Hz. A simetria em relação
    a $f_c$ confirma que o DSB-SC transmite informação redundante nas duas bandas
    laterais.

    \begin{nota}[Erros comuns]
      Medir incorretamente a posição dos picos por erros de cursor no espectro, ou não
      identificar a simetria das bandas laterais em torno de $f_c$.
    \end{nota}

  \item \textbf{O que acontece com o sinal demodulado quando há erro de fase no oscilador local?}

    \textbf{Resposta esperada:} A amplitude demodulada é proporcional a $\cos(\theta)$,
    onde $\theta$ é o erro de fase. Para $\theta=30°$: $\cos(30°) \approx 0{,}866$
    (atenuação de $13{,}4\%$); para $60°$: $\cos(60°) = 0{,}5$ (atenuação de $50\%$);
    para $90°$: $\cos(90°) = 0$ (cancelamento total). Este é o principal problema do
    DSB-SC: exige oscilador local perfeitamente sincronizado em fase.

    \begin{nota}[Erros comuns]
      Descrever apenas qualitativamente (``sinal cai'') sem calcular os valores de
      $\cos\theta$ para os erros de fase específicos indicados no roteiro.
    \end{nota}

  \item \textbf{Quais são os critérios para escolher a frequência de corte do filtro passa-baixas na demodulação?}

    \textbf{Resposta esperada:} O filtro deve ter frequência de corte entre $f_m$ e
    $2f_c - f_m$. Um corte muito baixo ($ < f_m$) atenua o próprio sinal mensagem.
    Um corte muito alto ($ > f_c$) permite que os componentes dobrados em $2f_c$
    passem e corrompam o sinal. O valor ótimo é $f_\text{cutoff} \approx 1{,}5\,f_m$.
    A largura de transição também é relevante para a atenuação dos termos em $2f_c$.

    \begin{nota}[Erros comuns]
      Não perceber que o LPF serve para eliminar os termos em $2f_c$ gerados pelo
      produto de demodulação, focando apenas na atenuação de ruído fora de banda.
    \end{nota}

  \item \textbf{Como a mensagem recuperada se compara com a mensagem original em amplitude e forma?}

    \textbf{Resposta esperada:} O sinal recuperado deve ser $\tfrac{A_m}{2}\cos(2\pi f_m t)$,
    ou seja, metade da amplitude original (fator $0{,}5$ do produto de dois cossenos).
    Esta atenuação é esperada e pode ser compensada por um amplificador de ganho 2.
    Se o LPF e o oscilador local estiverem corretamente configurados, não há distorção:
    apenas a escala difere da mensagem original.

    \begin{nota}[Erros comuns]
      Confundir o fator de escala esperado ($0{,}5$) com distorção ou erro de
      implementação do fluxograma.
    \end{nota}

\end{enumerate}

\section*{Critérios de correção}
\begin{tabular}{ll}
  \toprule
  Critério & Pontuação \\
  \midrule
  Q1 -- Ausência de portadora ligada à falta de DC em $m(t)$, com algebra & 1,0 ponto \\
  Q2 -- Cálculo correto de $f_c \pm f_m$ para cada valor de $f_m$         & 1,0 ponto \\
  Q3 -- Valores de $\cos\theta$ e impacto na amplitude demodulada          & 1,0 ponto \\
  Q4 -- Critério de escolha de $f_\text{cutoff}$ com justificativa         & 1,0 ponto \\
  Q5 -- Fator de escala 0,5 identificado e interpretado corretamente       & 1,0 ponto \\
  \midrule
  \textbf{Total} & \textbf{5,0 pontos} \\
  \bottomrule
\end{tabular}

\bigskip
\begin{nota}[Nota de correção geral]
  Priorizar coerência física e evidência experimental (capturas de tela, tabelas).
  Respostas sem evidência visual devem ter desconto. Não penalizar pequenos erros
  de cálculo se o raciocínio físico estiver correto.
\end{nota}

\end{document}
